\chapter{Parent-origin полной портальной матрицы двух источников}
\label{ch:v8-extra-mass-two-source-portal-origin}

\section{Прямая сумма не создаёт портал}

Пусть $X_{BM}$ несёт активные радиусы
$T_B=\Tr(BB^\dagger)$ и $T_M=\Tr(MM^\dagger)$, а
$h^0=aA+bB$ действует на отдельном центральном носителе. Для буквальной
прямой суммы
\begin{equation}
 X_\oplus=X_{BM}\oplus h^0
 \label{eq:v8-two-source-portal-direct-sum}
\end{equation}
любой полиномиальный момент аддитивен:
\begin{equation}
 \Tr X_\oplus^n=\Tr X_{BM}^n+\Tr(h^0)^n.
 \label{eq:v8-two-source-portal-additivity}
\end{equation}
Поэтому все смешанные производные по $(T_B,T_M)$ и $(a,b)$ равны нулю.
Для ненулевого портала необходим общий блоковый оператор, а не соседство
двух функционалов.

\section{Минимальный общий кубический блок}

Для одного скалярного ребра $F$ с $T=\Tr(FF^\dagger)$ и центральным
значением $\lambda$ рассмотрим
\begin{equation}
 X(\lambda,F)=\begin{pmatrix}\lambda I&F\\F^\dagger&0\end{pmatrix}.
 \label{eq:v8-two-source-portal-common-block}
\end{equation}
Прямое умножение даёт
\begin{equation}
 \Tr X^3=(\dim I)\lambda^3+3\lambda T.
 \label{eq:v8-two-source-portal-cubic-identity}
\end{equation}
Итак, кубический момент действительно создаёт линейный портал. Именно
конкретная матричная конечная геометрия определяет допустимые скалярные
связи spectral action \cite{v8-two-source-portal-krajewski}; формула сама
по себе не выбирает вложение $F$ в центральные сектора.

На блоках $(Y,u,d)$ направления $A,B$ имеют charge-векторы
\begin{equation}
 q_Y=(0,-3),\qquad q_u=(1,1),\qquad q_d=(-1,1),
 \label{eq:v8-two-source-portal-charge-vectors}
\end{equation}
то есть $\lambda_e=q_e\cdot(a,b)$. Назначив носителям $T_B,T_M$ сектора
$e_B,e_M\in\{Y,u,d\}$ и коэффициенты $c_B,c_M$, получаем
\begin{equation}
 V_3=3c_BT_B\,q_{e_B}\cdot(a,b)
     +3c_MT_M\,q_{e_M}\cdot(a,b)+V_{\rm self}^{(3)}.
 \label{eq:v8-two-source-portal-incidence-potential}
\end{equation}

\section{Ранговая классификация incidence-назначений}

Портальная матрица формы определяется строками
\begin{equation}
 Q(e_B,e_M)=\begin{pmatrix}q_{e_B}\\q_{e_M}\end{pmatrix}.
 \label{eq:v8-two-source-portal-incidence-matrix}
\end{equation}
Среди девяти упорядоченных назначений три диагональных имеют ранг один, а
шесть назначений в разные сектора имеют ранг два:
\begin{equation}
 N_{\rank Q=1}=3,qquad N_{\rank Q=2}=6,
 \qquad \det Q\in\{0,\pm2,\pm3\}.
 \label{eq:v8-two-source-portal-rank-menu}
\end{equation}
При двух независимых коэффициентах карта в источники равна
\begin{equation}
 \begin{pmatrix}j_A\\j_B\end{pmatrix}
 =-3\begin{pmatrix}T_Bq_{e_B}^{\,T}&T_Mq_{e_M}^{\,T}\end{pmatrix}
 \begin{pmatrix}c_B\\c_M\end{pmatrix},
 \qquad
 |\det|\in\{108,162\}
 \label{eq:v8-two-source-portal-two-coefficient-map}
\end{equation}
для $T_B=3,T_M=2$ и $e_B\ne e_M$. Следовательно, условный общий оператор
может реализовать полную rank-two source-map.

Один общий коэффициент $c_3$ оставляет только дискретное меню лучей:
\begin{equation}
 \frac1{3c_3}(j_A,j_B)
 \in\{(2,-7),(-2,-7),(3,-3),(1,5),(-3,-3),(-1,5)\}.
 \label{eq:v8-two-source-portal-single-coefficient-rays}
\end{equation}
Ни один такой однопараметрический луч не покрывает плоскость источников.

\section{Почему parent-origin остаётся открытым}

Текущий родитель содержит только чётные моменты общего спектрального
оператора. Поэтому
\begin{equation}
 c_3^{\rm inherited}=c_B^{\rm inherited}=c_M^{\rm inherited}=0.
 \label{eq:v8-two-source-portal-odd-coefficients-zero}
\end{equation}
Кроме того, активные матрицы $B,M$ и центральный носитель дополнительных
рёбер сейчас входят раздельными блоками. Для превращения
\eqref{eq:v8-two-source-portal-cubic-identity} в член действия нужны новые
типизированные incidence-вложения. Новые скалярные поля и связи в конечной
спектральной геометрии действительно возможны, но требуют изменения её
Dirac-блоков \cite{v8-two-source-portal-stephan}.

Шесть rank-two назначений не эквивалентны: гиперзаряды и цвет различают
$Y,u,d$, а сохранённая симметрия не переставляет эти сектора. Поэтому
никакое из них не выбирается текущей группой. Реестры равны
\begin{equation}
 N_{\rm conditional\ operator\ shape}=7/7,
 \qquad N_{\rm incidence\ selector}=0/5,
 \qquad N_{\rm parent\ origin}=0/4.
 \label{eq:v8-two-source-portal-ledgers}
\end{equation}

\begin{theorem}[Условная rank-two кубическая реализация]
Минимальный общий блок удовлетворяет $\Tr X^3=d\lambda^3+3\lambda T$.
Если два ненулевых радиуса назначены двум различным центральным секторам и
имеют независимые кубические коэффициенты, карта коэффициентов в источники
имеет ненулевой определитель $108$ или $162$. Поэтому полная портальная
форма существует условно.
\end{theorem}

\begin{nogocriterion}[Меню incidence не является parent-origin]
Нельзя объявлять полную portal-matrix выведенной: прямой сумме недостаёт
смешанных моментов, текущий parent не содержит нечётных коэффициентов, а
шесть rank-two incidence-назначений дают разные физические лучи. Выбор
назначения после просмотра желаемых щелей является ручным селектором.
\end{nogocriterion}

\begin{protocol}\begin{enumerate}
\item прямая сумма создаёт портал: \textbf{нет};
\item минимальный смешанный момент: \textbf{кубический};
\item коэффициент $\lambda T$ в $\Tr X^3$: \textbf{$3$};
\item центральных charge-векторов: \textbf{три};
\item упорядоченных incidence-назначений: \textbf{девять};
\item rank-two назначений: \textbf{шесть};
\item два независимых кубических коэффициента достаточны: \textbf{да условно};
\item один общий коэффициент покрывает плоскость: \textbf{нет};
\item inherited odd coefficients: \textbf{ноль};
\item conditional shape: \textbf{$7/7$};
\item incidence selector: \textbf{$0/5$};
\item parent-origin: \textbf{$0/4$};
\item следующий гейт: \textbf{селектор кубического incidence-назначения}.
\end{enumerate}\end{protocol}

Машинный сертификат сохранён в
\path{s2t_v8_baryon_c0_singlet_triplet_central_gap_isotypic_channel_extra_edge_mass_two_source_portal_matrix_parent_origin_gate_results.json}.

\begin{thebibliography}{9}
\bibitem{v8-two-source-portal-krajewski} T.~Krajewski,
\emph{Constraints on the Scalar Potential from Spectral Action Principle},
arXiv:hep-th/9803199.
\bibitem{v8-two-source-portal-stephan} C.~A.~Stephan,
\emph{New Scalar Fields in Noncommutative Geometry}, arXiv:0901.4676.
\end{thebibliography}